\documentclass[11pt,a4paper]{article}

% ============================================================
% PACKAGES
% ============================================================
\usepackage[utf8]{inputenc}
\usepackage[T1]{fontenc}
\usepackage[french]{babel}
\usepackage{lmodern}
\usepackage{geometry}
\geometry{margin=1.6cm}
\usepackage{xcolor}
\usepackage{amsmath,amssymb}
\usepackage{tikz}
\usepackage{array}
\usepackage{enumitem}
\usepackage{multicol}
\usepackage{fancyhdr}
\usepackage{tcolorbox}

% ============================================================
% STYLE
% ============================================================
\pagestyle{fancy}
\fancyhf{}
\lhead{\textbf{Mathématiques -- Sixième}}
\rhead{\textbf{Fiche méthode -- Symétrie axiale}}
\cfoot{\thepage}

\setlength{\parindent}{0pt}
\setlength{\parskip}{0.35em}

\definecolor{bleu}{RGB}{30,80,160}
\definecolor{rouge}{RGB}{180,40,40}
\definecolor{vert}{RGB}{20,130,70}
\definecolor{grisclair}{RGB}{245,245,245}
\definecolor{orange}{RGB}{220,120,20}

\newcommand{\important}[1]{\textcolor{rouge}{\textbf{#1}}}
\newcommand{\methode}[1]{\textcolor{bleu}{\textbf{#1}}}

\newtcolorbox{boite}[1]{
  colback=grisclair,
  colframe=bleu,
  title=\textbf{#1},
  fonttitle=\bfseries
}

\newtcolorbox{attention}{
  colback=red!5,
  colframe=rouge,
  title=\textbf{Attention},
  fonttitle=\bfseries
}

\newtcolorbox{retenir}{
  colback=green!5,
  colframe=vert,
  title=\textbf{À retenir},
  fonttitle=\bfseries
}

\begin{document}

% ============================================================
% TITRE
% ============================================================
\begin{center}
    {\Large \textbf{Fiche méthode}}\\[0.15cm]
    {\LARGE \textbf{Tracer le symétrique d'un point ou d'une figure}}\\[0.15cm]
    {\large Niveau sixième -- Symétrie axiale}
\end{center}

\vspace{0.3cm}

\begin{boite}{Objectif de la fiche}
Cette fiche explique les différentes méthodes pour construire le symétrique d'un point ou d'une figure par rapport à une droite appelée \textbf{axe de symétrie}.  

On verra les constructions avec :
\begin{itemize}[label=\textbullet]
    \item la règle graduée et l'équerre ;
    \item la règle non graduée et l'équerre ;
    \item la règle et le compas ;
    \item le compas seul ;
    \item les cas particuliers : point sur l'axe, figure qui coupe l'axe, figure complète à compléter.
\end{itemize}
\end{boite}

% ============================================================
\section*{1. Définition : symétrique d'un point}

Soit une droite $(d)$ et un point $A$.  
Le point $A'$ est le symétrique de $A$ par rapport à la droite $(d)$ lorsque la droite $(d)$ est la \textbf{médiatrice} du segment $[AA']$.

\begin{retenir}
Dire que $(d)$ est la médiatrice de $[AA']$ signifie deux choses :
\begin{enumerate}[label=\arabic*)]
    \item $(d)$ est \important{perpendiculaire} au segment $[AA']$ ;
    \item $(d)$ passe par le \important{milieu} du segment $[AA']$.
\end{enumerate}
\end{retenir}

\begin{center}
\begin{tikzpicture}[scale=1]

% axe
\draw[very thick, bleu] (0,-2) -- (0,2);
\node[bleu, above] at (0,2) {$(d)$};

% points
\filldraw (-2,0.8) circle (2pt) node[left] {$A$};
\filldraw (2,0.8) circle (2pt) node[right] {$A'$};

% segment
\draw[thick, dashed] (-2,0.8)--(2,0.8);

% angle droit
\draw (0,0.8) ++(0,0) rectangle ++(0.25,0.25);

% milieu
\filldraw (0,0.8) circle (1.5pt) node[below right] {$I$};

% égalités
\draw[<->] (-2,0.45)--(0,0.45);
\draw[<->] (0,0.45)--(2,0.45);
\node at (-1,0.25) {$AI$};
\node at (1,0.25) {$IA'$};
\node at (0,-1.7) {$(d)$ est perpendiculaire à $[AA']$ et passe par son milieu.};

\end{tikzpicture}
\end{center}

% ============================================================
\section*{2. Méthode 1 : avec règle graduée et équerre}

\begin{boite}{Situation}
On veut construire le symétrique d'un point $A$ par rapport à une droite $(d)$ en utilisant une \textbf{règle graduée} et une \textbf{équerre}.
\end{boite}

\subsection*{Étapes}

\begin{enumerate}[label=\textbf{Étape \arabic* :}]
    \item Placer l'équerre pour tracer la droite perpendiculaire à $(d)$ passant par $A$.
    \item Cette perpendiculaire coupe $(d)$ en un point $I$.
    \item Mesurer la longueur $AI$ avec la règle graduée.
    \item De l'autre côté de $(d)$, placer le point $A'$ sur la même droite, de façon que :
    \[
    IA'=IA.
    \]
    \item Le point $A'$ est le symétrique de $A$ par rapport à $(d)$.
\end{enumerate}

\begin{center}
\begin{tikzpicture}[scale=1]

% axe
\draw[very thick, bleu] (0,-2) -- (0,2);
\node[bleu, above] at (0,2) {$(d)$};

% point A
\filldraw (-2,1) circle (2pt) node[left] {$A$};

% perpendiculaire
\draw[dashed, thick, rouge] (-2,1)--(2.2,1);

% intersection
\filldraw (0,1) circle (1.5pt) node[below right] {$I$};

% point A'
\filldraw (2,1) circle (2pt) node[right] {$A'$};

% angle droit
\draw (0,1) rectangle +(0.25,0.25);

% flèches distances
\draw[<->, vert, thick] (-2,0.65)--(0,0.65);
\draw[<->, vert, thick] (0,0.65)--(2,0.65);
\node[vert] at (-1,0.45) {$AI$};
\node[vert] at (1,0.45) {$IA'$};

\node at (0,-1.6) {On trace la perpendiculaire puis on reporte la même distance.};

\end{tikzpicture}
\end{center}

\begin{retenir}
Avec la règle graduée et l'équerre, on utilise directement la propriété :
\[
(d) \perp [AA']
\quad \text{et} \quad
AI=IA'.
\]
\end{retenir}

% ============================================================
\section*{3. Méthode 2 : avec règle non graduée et équerre}

\begin{boite}{Situation}
On veut construire le symétrique d'un point $A$ par rapport à une droite $(d)$ avec une \textbf{règle non graduée} et une \textbf{équerre}.  
La règle sert seulement à tracer des droites, pas à mesurer.
\end{boite}

\subsection*{Étapes}

\begin{enumerate}[label=\textbf{Étape \arabic* :}]
    \item Avec l'équerre, tracer la perpendiculaire à $(d)$ passant par $A$.
    \item Elle coupe $(d)$ en un point $I$.
    \item Pour placer $A'$, on doit reporter la distance $AI$ sans utiliser les graduations.
    \item On peut utiliser un compas pour reporter la distance, ou utiliser une construction déjà imposée par l'enseignant.
\end{enumerate}

\begin{attention}
Avec une règle non graduée et une équerre seulement, on peut tracer la perpendiculaire, mais on ne peut pas mesurer précisément la distance $AI$.  
Pour une construction exacte, il faut souvent ajouter le compas.
\end{attention}

\begin{center}
\begin{tikzpicture}[scale=1]

% axe oblique
\draw[very thick, bleu] (-2,-1.4)--(2,1.4);
\node[bleu, right] at (2,1.4) {$(d)$};

% point A
\filldraw (-2,1.3) circle (2pt) node[left] {$A$};

% projection approx
\coordinate (I) at (-0.5,0);
\coordinate (Ap) at (1,-1.3);

% ligne perpendiculaire
\draw[dashed, thick, rouge] (-2,1.3)--(1,-1.3);

% points
\filldraw (I) circle (1.5pt) node[below right] {$I$};
\filldraw (Ap) circle (2pt) node[right] {$A'$};

\node at (0,-2.2) {L'équerre sert à tracer la perpendiculaire à l'axe.};

\end{tikzpicture}
\end{center}

% ============================================================
\section*{4. Méthode 3 : avec règle et compas}

\begin{boite}{Situation}
On veut construire le symétrique de $A$ par rapport à $(d)$ sans mesurer avec une règle graduée.  
On utilise une \textbf{règle non graduée} et un \textbf{compas}.
\end{boite}

\subsection*{Idée importante}

On choisit deux points $M$ et $N$ sur l'axe de symétrie $(d)$.  
Si $A'$ est le symétrique de $A$, alors :
\[
MA=MA'
\qquad \text{et} \qquad
NA=NA'.
\]

Ainsi, $A'$ est à l'intersection :
\begin{itemize}[label=\textbullet]
    \item du cercle de centre $M$ passant par $A$ ;
    \item du cercle de centre $N$ passant par $A$.
\end{itemize}

\subsection*{Étapes}

\begin{enumerate}[label=\textbf{Étape \arabic* :}]
    \item Choisir deux points $M$ et $N$ sur la droite $(d)$.
    \item Ouvrir le compas avec l'écartement $MA$.
    \item Tracer un arc de cercle de centre $M$ passant par $A$.
    \item Ouvrir le compas avec l'écartement $NA$.
    \item Tracer un arc de cercle de centre $N$ passant par $A$.
    \item Les deux arcs se coupent en $A$ et en un autre point : ce deuxième point est $A'$.
\end{enumerate}

\begin{center}
\begin{tikzpicture}[scale=0.95]

% axe
\draw[very thick, bleu] (-4,0)--(4,0);
\node[bleu, right] at (4,0) {$(d)$};

% points M N
\filldraw[bleu] (-2,0) circle (2pt) node[below] {$M$};
\filldraw[bleu] (2,0) circle (2pt) node[below] {$N$};

% A and A'
\filldraw (0,2) circle (2pt) node[above] {$A$};
\filldraw (0,-2) circle (2pt) node[below] {$A'$};

% arcs/circles
\draw[orange, thick] (-2,0) circle ({sqrt(8)});
\draw[vert, thick] (2,0) circle ({sqrt(8)});

% dashed vertical
\draw[dashed, rouge] (0,2)--(0,-2);

\node[orange] at (-3.4,2.1) {cercle de centre $M$};
\node[vert] at (3.2,2.1) {cercle de centre $N$};

\end{tikzpicture}
\end{center}

\begin{retenir}
Cette méthode est très utile lorsque l'axe est oblique, car elle évite de mesurer avec une règle graduée.
\end{retenir}

% ============================================================
\section*{5. Méthode 4 : avec le compas seul}

\begin{boite}{Situation}
On veut construire le symétrique d'un point $A$ par rapport à une droite $(d)$ en utilisant \textbf{seulement le compas}.  
On suppose que deux points $M$ et $N$ de l'axe $(d)$ sont déjà donnés.
\end{boite}

\subsection*{Principe}

On ne trace pas forcément la droite, on utilise seulement les distances.

Si $A'$ est le symétrique de $A$ par rapport à $(d)$, alors :
\[
MA=MA'
\qquad \text{et} \qquad
NA=NA'.
\]

Donc $A'$ est l'autre intersection des deux cercles :
\[
\mathcal{C}(M,MA)
\qquad \text{et} \qquad
\mathcal{C}(N,NA).
\]

\subsection*{Étapes}

\begin{enumerate}[label=\textbf{Étape \arabic* :}]
    \item Piquer le compas en $M$ et prendre l'écartement $MA$.
    \item Tracer un arc de cercle passant par $A$.
    \item Piquer le compas en $N$ et prendre l'écartement $NA$.
    \item Tracer un second arc de cercle passant par $A$.
    \item Les deux arcs se coupent en $A$ et en un autre point $A'$.
    \item Le point $A'$ est le symétrique de $A$ par rapport à la droite $(MN)$.
\end{enumerate}

\begin{center}
\begin{tikzpicture}[scale=1]

% axe oblique
\draw[very thick, bleu] (-3,-1)--(3,1);
\node[bleu, right] at (3,1) {$(d)$};

% points M N on axis
\filldraw[bleu] (-1.5,-0.5) circle (2pt) node[below left] {$M$};
\filldraw[bleu] (1.5,0.5) circle (2pt) node[above right] {$N$};

% point A and approximate reflected A'
\filldraw (-2,2) circle (2pt) node[above left] {$A$};
\filldraw (-0.9,-1.3) circle (2pt) node[below] {$A'$};

% arcs approximate using circles
\draw[orange, thick] (-1.5,-0.5) circle (2.55);
\draw[vert, thick] (1.5,0.5) circle (3.82);

\node[orange] at (-3.7,0.2) {rayon $MA$};
\node[vert] at (3.7,-0.9) {rayon $NA$};

\end{tikzpicture}
\end{center}

\begin{attention}
Avec le compas seul, il faut deux points $M$ et $N$ sur l'axe.  
Sans ces deux points, on ne peut pas utiliser cette méthode directement.
\end{attention}

% ============================================================
\section*{6. Tracer le symétrique d'un segment}

\begin{boite}{Situation}
On veut construire le symétrique du segment $[AB]$ par rapport à une droite $(d)$.
\end{boite}

\subsection*{Méthode}

\begin{enumerate}[label=\textbf{Étape \arabic* :}]
    \item Construire le symétrique $A'$ du point $A$.
    \item Construire le symétrique $B'$ du point $B$.
    \item Tracer le segment $[A'B']$.
\end{enumerate}

\begin{center}
\begin{tikzpicture}[scale=0.9]

% axis
\draw[very thick, bleu] (0,-2.5)--(0,2.5);
\node[bleu, above] at (0,2.5) {$(d)$};

% original segment
\filldraw (-3,1.5) circle (2pt) node[left] {$A$};
\filldraw (-2,-1) circle (2pt) node[left] {$B$};
\draw[very thick] (-3,1.5)--(-2,-1);

% reflected segment
\filldraw (3,1.5) circle (2pt) node[right] {$A'$};
\filldraw (2,-1) circle (2pt) node[right] {$B'$};
\draw[very thick, rouge] (3,1.5)--(2,-1);

% construction lines
\draw[dashed] (-3,1.5)--(3,1.5);
\draw[dashed] (-2,-1)--(2,-1);

\end{tikzpicture}
\end{center}

\begin{retenir}
Le symétrique d'un segment est un segment de même longueur.
\[
AB=A'B'.
\]
\end{retenir}

% ============================================================
\section*{7. Tracer le symétrique d'un triangle}

\begin{boite}{Situation}
On veut construire le symétrique d'un triangle $ABC$ par rapport à une droite $(d)$.
\end{boite}

\subsection*{Méthode}

\begin{enumerate}[label=\textbf{Étape \arabic* :}]
    \item Construire le symétrique $A'$ du sommet $A$.
    \item Construire le symétrique $B'$ du sommet $B$.
    \item Construire le symétrique $C'$ du sommet $C$.
    \item Relier $A'$, $B'$ et $C'$.
\end{enumerate}

\begin{center}
\begin{tikzpicture}[scale=0.9]

% axis oblique
\draw[very thick, bleu] (-1,-2.5)--(1,2.5);
\node[bleu, above] at (1,2.5) {$(d)$};

% original triangle
\coordinate (A) at (-4,1.8);
\coordinate (B) at (-2,2.2);
\coordinate (C) at (-3,-0.8);

\draw[very thick] (A)--(B)--(C)--cycle;
\filldraw (A) circle (2pt) node[left] {$A$};
\filldraw (B) circle (2pt) node[above] {$B$};
\filldraw (C) circle (2pt) node[left] {$C$};

% approximate reflected triangle
\coordinate (Ap) at (3.1,0.1);
\coordinate (Bp) at (2.2,-1.7);
\coordinate (Cp) at (0.5,0.2);

\draw[very thick, rouge] (Ap)--(Bp)--(Cp)--cycle;
\filldraw[rouge] (Ap) circle (2pt) node[right] {$A'$};
\filldraw[rouge] (Bp) circle (2pt) node[right] {$B'$};
\filldraw[rouge] (Cp) circle (2pt) node[right] {$C'$};

\node at (0,-3) {On construit l'image de chaque sommet, puis on relie.};

\end{tikzpicture}
\end{center}

\begin{retenir}
La symétrie axiale conserve :
\begin{itemize}[label=\textbullet]
    \item les longueurs ;
    \item les angles ;
    \item les formes ;
    \item les alignements.
\end{itemize}
\end{retenir}

% ============================================================
\section*{8. Tracer le symétrique d'un polygone}

\begin{boite}{Situation}
On veut construire le symétrique d'un polygone $ABCDE$ par rapport à une droite $(d)$.
\end{boite}

\subsection*{Méthode}

\begin{enumerate}[label=\textbf{Étape \arabic* :}]
    \item Construire le symétrique de chaque sommet :
    \[
    A',B',C',D',E'.
    \]
    \item Relier les points dans le même ordre que dans la figure de départ.
    \item Ne pas oublier de fermer la figure.
\end{enumerate}

\begin{attention}
Il faut relier les points dans le même ordre.  
Si la figure de départ est $A-B-C-D-E$, alors la figure symétrique est $A'-B'-C'-D'-E'$.
\end{attention}

\begin{center}
\begin{tikzpicture}[scale=0.85]

% axis
\draw[very thick, bleu] (0,-2.5)--(0,3);
\node[bleu, above] at (0,3) {$(d)$};

% original polygon
\coordinate (A) at (-4,2);
\coordinate (B) at (-2.5,2.5);
\coordinate (C) at (-1.7,0.8);
\coordinate (D) at (-2.5,-1.4);
\coordinate (E) at (-4,-0.6);

\draw[very thick] (A)--(B)--(C)--(D)--(E)--cycle;
\foreach \p/\name in {A/A,B/B,C/C,D/D,E/E}{
  \filldraw (\p) circle (2pt) node[left] {$\name$};
}

% reflected polygon
\coordinate (Ap) at (4,2);
\coordinate (Bp) at (2.5,2.5);
\coordinate (Cp) at (1.7,0.8);
\coordinate (Dp) at (2.5,-1.4);
\coordinate (Ep) at (4,-0.6);

\draw[very thick, rouge] (Ap)--(Bp)--(Cp)--(Dp)--(Ep)--cycle;
\foreach \p/\name in {Ap/A',Bp/B',Cp/C',Dp/D',Ep/E'}{
  \filldraw[rouge] (\p) circle (2pt) node[right] {$\name$};
}

\end{tikzpicture}
\end{center}

% ============================================================
\section*{9. Cas particulier : un point est sur l'axe}

\begin{boite}{Situation}
Le point $A$ appartient à l'axe de symétrie $(d)$.
\end{boite}

Dans ce cas, le symétrique de $A$ est $A$ lui-même.

\[
A'=A.
\]

\begin{center}
\begin{tikzpicture}[scale=1]

% axis
\draw[very thick, bleu] (-3,0)--(3,0);
\node[bleu, right] at (3,0) {$(d)$};

% point on axis
\filldraw (0,0) circle (2pt) node[above] {$A=A'$};

\node at (0,-1) {Un point placé sur l'axe ne bouge pas.};

\end{tikzpicture}
\end{center}

\begin{retenir}
Tout point situé sur l'axe de symétrie est son propre symétrique.
\end{retenir}

% ============================================================
\section*{10. Cas particulier : la figure coupe l'axe}

\begin{boite}{Situation}
Une figure coupe l'axe de symétrie. Certains points sont donc sur l'axe.
\end{boite}

\subsection*{Méthode}

\begin{itemize}[label=\textbullet]
    \item Les points situés sur l'axe ne bougent pas.
    \item Les autres points sont construits normalement.
    \item On relie ensuite les points dans le bon ordre.
\end{itemize}

\begin{center}
\begin{tikzpicture}[scale=0.9]

% axis
\draw[very thick, bleu] (0,-2.5)--(0,2.5);
\node[bleu, above] at (0,2.5) {$(d)$};

% original half figure
\coordinate (A) at (0,2);
\coordinate (B) at (-2,1);
\coordinate (C) at (-1,-1);
\coordinate (D) at (0,-2);

\draw[very thick] (A)--(B)--(C)--(D);
\filldraw (A) circle (2pt) node[right] {$A=A'$};
\filldraw (B) circle (2pt) node[left] {$B$};
\filldraw (C) circle (2pt) node[left] {$C$};
\filldraw (D) circle (2pt) node[right] {$D=D'$};

% reflected
\coordinate (Bp) at (2,1);
\coordinate (Cp) at (1,-1);

\draw[very thick, rouge] (A)--(Bp)--(Cp)--(D);
\filldraw[rouge] (Bp) circle (2pt) node[right] {$B'$};
\filldraw[rouge] (Cp) circle (2pt) node[right] {$C'$};

\end{tikzpicture}
\end{center}

% ============================================================
\section*{11. Compléter une figure par symétrie}

\begin{boite}{Situation}
On te donne une moitié de figure et un axe de symétrie.  
Il faut compléter l'autre moitié.
\end{boite}

\subsection*{Méthode}

\begin{enumerate}[label=\textbf{Étape \arabic* :}]
    \item Repérer tous les points importants de la figure.
    \item Pour chaque point qui n'est pas sur l'axe, construire son symétrique.
    \item Les points situés sur l'axe restent fixes.
    \item Relier les points obtenus en suivant la même forme que la figure de départ.
\end{enumerate}

\begin{center}
\begin{tikzpicture}[scale=0.9]

% axis
\draw[very thick, bleu] (0,-2.8)--(0,2.8);
\node[bleu, above] at (0,2.8) {$(d)$};

% left half
\draw[very thick] (0,2)--(-2,1)--(-2,-1)--(0,-2);
\filldraw (0,2) circle (2pt) node[right] {$A$};
\filldraw (-2,1) circle (2pt) node[left] {$B$};
\filldraw (-2,-1) circle (2pt) node[left] {$C$};
\filldraw (0,-2) circle (2pt) node[right] {$D$};

% right half
\draw[very thick, rouge] (0,2)--(2,1)--(2,-1)--(0,-2);
\filldraw[rouge] (2,1) circle (2pt) node[right] {$B'$};
\filldraw[rouge] (2,-1) circle (2pt) node[right] {$C'$};

\node at (0,-3.4) {Les points sur l'axe restent en place.};

\end{tikzpicture}
\end{center}

% ============================================================
\section*{12. Choisir la bonne méthode selon les instruments}

\begin{center}
\renewcommand{\arraystretch}{1.4}
\begin{tabular}{|p{0.28\linewidth}|p{0.36\linewidth}|p{0.28\linewidth}|}
\hline
\textbf{Instruments autorisés} & \textbf{Méthode} & \textbf{Idée principale} \\
\hline
Règle graduée + équerre 
& Tracer la perpendiculaire, mesurer, reporter 
& Utiliser le milieu et la perpendicularité \\
\hline
Règle non graduée + équerre 
& Tracer la perpendiculaire 
& Utile pour obtenir la direction du segment $[AA']$ \\
\hline
Règle + compas 
& Utiliser deux points $M$ et $N$ sur l'axe 
& Le symétrique conserve les distances à $M$ et $N$ \\
\hline
Compas seul 
& Tracer deux arcs de cercle de centres $M$ et $N$ 
& $A'$ est l'autre intersection des deux arcs \\
\hline
Papier quadrillé 
& Compter les carreaux perpendiculairement à l'axe 
& Reporter le même nombre de carreaux de l'autre côté \\
\hline
\end{tabular}
\end{center}

% ============================================================
\section*{13. Erreurs classiques à éviter}

\begin{attention}
\begin{enumerate}[label=\arabic*)]
    \item Ne pas placer le point symétrique au hasard de l'autre côté de l'axe.
    \item Ne pas oublier que le segment $[AA']$ doit être perpendiculaire à l'axe.
    \item Ne pas oublier que l'axe doit passer par le milieu de $[AA']$.
    \item Ne pas déplacer les points qui sont déjà sur l'axe.
    \item Ne pas relier les sommets dans le mauvais ordre.
    \item Ne pas effacer tous les traits de construction : ils permettent de vérifier la méthode.
\end{enumerate}
\end{attention}

% ============================================================
\section*{14. Résumé général}

\begin{retenir}
Pour construire le symétrique d'une figure par rapport à une droite :
\begin{enumerate}[label=\arabic*)]
    \item On construit le symétrique de chaque point important.
    \item On vérifie que l'axe est la médiatrice entre chaque point et son image.
    \item On relie les points obtenus dans le même ordre.
\end{enumerate}
\end{retenir}

\[
\boxed{
\text{La symétrie axiale conserve les longueurs, les angles, les formes et les alignements.}
}
\]

\vspace{0.5cm}

\begin{center}
\textbf{Fin de la fiche méthode}
\end{center}

\end{document}